\section{OBJECTIVE AND SOURCE REVIEW}
\textbf{Erd\H{o}s \#138; independent attempt, 2026-09-23.}
The main target is $W(k)^{1/k}\to\infty$ for two colours. I read the entire supplied \texttt{138.tex}. Its Berlekamp consequence is an elementary implication of a known theorem; its assertion that all overlap edges are forced in every dependency graph is false. This attempt corrects that specific obstruction claim.
\section{WORK: EXACT DEPENDENCE CALCULATION}
Let $P,Q$ be distinct $k$-term progressions and $E_P,E_Q$ their monochromatic events under independent unbiased colouring. Put $t=|P\cap Q|$. Each event has probability $2^{1-k}$. For $t\ge1$, their simultaneous occurrence forces every point of $P\cup Q$ to have one common colour, giving
\[
\mathbb P(E_P\cap E_Q)=2^{1-2k+t}
=2^{t-1}\mathbb P(E_P)\mathbb P(E_Q).
\]
Thus one-point intersections give \emph{independent} events. For example, $P=\{1,2,3\}$ and $Q=\{3,4,5\}$ have joint probability $1/16=(1/4)^2$. Intersections of at least two points are dependent.
The supplied degree argument remains valid for the \emph{chosen full overlap graph}: if its maximum degree is $D$, the sufficient test $e2^{1-k}(D+1)\le1$ cannot hold beyond its stated exponential range. However, failure of this sufficient test does not establish failure of the symmetric local lemma in every permissible graph, still less all local-lemma methods.
A further subtlety prevents simply deleting every one-point overlap edge: a dependency graph requires independence from the sigma-algebra generated by all non-neighbours, rather than only pairwise independence. The displayed calculation licenses no such graph deletion by itself.
\section{ADVERSARIAL VERIFICATION AND REMAINING GAP}
The cases $t=0$ and $t=1$ both give independence but for different reasons; $t=0$ uses disjoint coordinates. The algebra concerns two events only. No asymptotic lower bound for $W(k)$ has been improved here. In particular, an unbounded ratio $W(k)/2^k$ along prime-indexed subsequences does not imply a divergent $k$th root. The official entry also distinguishes the now-resolved consecutive-difference variant from the remaining two-colour root-growth problem.
\section{SOURCES CHECKED}
Complete local source: \texttt{gpt\_pro\_5.4/138.tex}. Official indexed entry checked 2026-09-23: \url{https://www.erdosproblems.com/138}. The probability correction is proved directly above.
\section{FINAL}
\textbf{UNRESOLVED.} A false dependency assertion is removed; no new root-growth bound is proved.
\section{COMPLETION ESTIMATE}
Subjective progress toward the full two-colour question.
\textbf{COMPLETION: 5\%}
